\documentclass[10pt, a4paper]{article}

% Parametri che modificano il file main.tex
% Le uniche parti da cambiare su main.tex sono:
% - tabella versioni
% - testo file

\def\titolo{Analisi dei \\ \ requisiti} % \\ per andare a capo, \ per spazio
\def\titoloHeader{Analisi dei requisiti}
\def\data{v2.0.0(0)}

% \def\listaComponenti{
% Bresolin G.,
% Campese M.,
% Ciriolo I.,
% Dugo A.,
% Feltrin E.,
% Michelon R.,
% Orlandi G.
% }

\usepackage{style}
\usepackage{headerfooter}

\title{\titolo}
\author{SWEetCode}

\begin{document}

% PRIMA PAGINA
\begin{titlepage}
    \thispagestyle{empty}
    \begin{tikzpicture}[remember picture, overlay]
        % TRIANGOLI
        \draw[fill=secondarycolor, secondarycolor] (current page.north west) -- (current page.south west) -- (8.8, -28);
        \draw[fill=primarycolor, primarycolor] (-3, 5) -- (4, -13.6) -- (11, 5);

        % LOGO
        \node [xshift=-5cm, yshift=25cm] (logo) at (current page.south east) {\includegraphics[width=6.5cm]{img/logo.png}};

        % SWEETCODE - DATE
        \node [anchor=north east, align=right, xshift=-1.2cm, yshift=20.5cm, text=black] (sweetcode) at (current page.south east) {\fontsize{32pt}{36pt}\selectfont SWEetCode};
        \draw[line width=4pt, lightcol] ([xshift=-3cm, yshift=-0.37cm]sweetcode.south west) -- ([yshift=-0.37cm]sweetcode.south east);
        \node [anchor=north east, align=right, xshift=-1.2cm, yshift=18.7cm, text=black] (date) at (current page.south east){\fontsize{24pt}{24pt} \selectfont Verbale \tipoVerb};

        % NOME FILE
        \node [anchor=north east, text width=15cm, align=right, xshift=-1.2cm, yshift=17cm, text=black] (titolo) at (current page.south east){\fontsize{48pt}{48pt}\textbf{\data}};

        % BOX DATI PARTECIPANTI
        \ifthenelse{\equal{\tipoVerb}{Esterno}}{
            \node[anchor=north east, xshift=-1.2cm, yshift=14.5cm, minimum width=8cm] (box) at (current page.south east){};
        }{
            \node[anchor=north east, xshift=-1.2cm, yshift=12.5cm, minimum width=8cm] (box) at (current page.south east){};
        }

        % RESPONSABILE
        \node[anchor=north west, align=left] (dati2) at (box.north west) {\fontsize{15pt}{15pt}\selectfont \textbf{Responsabile}};
        \draw[line width=4pt, lightcol] (dati2.south west) -- ([xshift=8cm]dati2.south west);
        \node[anchor=north west, align=left] (dati21) at (dati2.south west){\fontsize{13pt}{13pt}\selectfont \nomeResp};

        % VERIFICATORE
        \node[anchor=north west, yshift=-1cm, align=left] (dati3) at (dati21.north west) {\fontsize{15pt}{15pt}\selectfont \textbf{Verificatore}};
        \draw[line width=4pt, lightcol] (dati3.south west) -- ([xshift=8cm]dati3.south west);
        \node[anchor=north west, align=left] (dati31) at (dati3.south west){\fontsize{13pt}{13pt}\selectfont \nomeVer};

        % SEGRETARIO DI RIUNIONE
        \node[anchor=north west, yshift=-1cm, align=left] (dati4) at (dati31.north west) {\fontsize{15pt}{15pt}\selectfont \textbf{Segretario di Riunione}};
        \draw[line width=4pt, lightcol] (dati4.south west) -- ([xshift=8cm]dati4.south west);
        \node[anchor=north west, align=left] (dati41) at (dati4.south west){\fontsize{13pt}{13pt}\selectfont \nomeSegr};
        
        % UNIPD - SWE
        \node [xshift=4.4cm, yshift=2.3cm, draw, secondarycolor, text=white] (uni) at (current page.south west) {\fontsize{20pt}{20pt} \selectfont Università di Padova};
        \node [xshift=0.65cm, yshift=0.7cm, draw, secondarycolor, text=white, below=of uni] (corso) {\fontsize{20pt}{20pt}\selectfont Ingegneria del Software};

        % FIRMA AZIENDA
        \ifthenelse{\equal{\tipoVerb}{Esterno}}{
            \draw[line width=4pt, lightcol] ([xshift=-1.2cm, yshift=7cm]current page.south east) -- ([xshift=-8cm, yshift=7cm]current page.south east);
            \node [xshift=-4.8cm, yshift=8cm] (logo) at (current page.south east) {\includegraphics[width=6cm]{img/firme/\firmaAzienda}};
            \node[anchor=north west, xshift=12.9cm, yshift=6.65cm, align=left] at (current page.south west)
            {\fontsize{13pt}{13pt}\selectfont L'azienda: \nomeAzienda};
        }

        % FIRMA CODING COWBOYS
        \draw[line width=4pt, lightcol] ([xshift=-1.2cm, yshift=4.35cm]current page.south east) -- ([xshift=-8cm, yshift=4.35cm]current page.south east);
        \node [xshift=-4.8cm, yshift=5.5cm] (logo) at (current page.south east) {\includegraphics[width=6cm]{img/firme/menon.png}};
        \node[anchor=north west, xshift=12.9cm, yshift=4cm, align=left] at (current page.south west)
        {\fontsize{13pt}{13pt}\selectfont Coding Cowboys: Menon G.};

        % FIRMA
        \draw[line width=4pt, lightcol] ([xshift=-1.2cm, yshift=1.8cm]current page.south east) -- ([xshift=-8cm, yshift=1.8cm]current page.south east);
        \node [xshift=-4.8cm, yshift=2.45cm] (logo) at (current page.south east) {\includegraphics[width=6cm]{img/firme/\firmaResp}};
        \node[anchor=north west, xshift=12.9cm, yshift=1.45cm, align=left] at (current page.south west)
        {\fontsize{13pt}{13pt}\selectfont Il Responsabile: \nomeResp};
        
    \end{tikzpicture}
\end{titlepage}

% REGISTRO DELLE VERSIONI
%Registro in ordine dalla più recente alla meno recente!
{\renewcommand{\arraystretch}{1.5}
\section*{Registro delle versioni}

\begin{xltabular}{\textwidth}{c|c|c|c|X}
\label{tab:long}

\textbf{Versione} & \textbf{Data} & \quantities{\textbf{Responsabile di}\\\textbf{stesura}}& \textbf{Revisore} & \quantities{\textbf{Dettaglio e}\\\textbf{motivazioni}} \\
\endfirsthead

\textbf{Versione} & \textbf{Data} & \quantities{\textbf{Responsabile di}\\\textbf{stesura}}& \textbf{Revisore} & \quantities{\textbf{Dettaglio e}\\\textbf{motivazioni}} \\
\endhead

\multicolumn{5}{r}{{Continua nella pagina successiva}} \\
\endfoot

\endlastfoot

\hline
v1.10.0(7) & $2023-12-29$ & \quantities{Michelon R.} & Dugo A. & Aggiunti diagrammi use case e di attività.\\
\hline
v1.10.0(6) & $2023-12-29$ & \quantities{Michelon R.} & Dugo A. & Aggiornamento use case e requisiti.\\
\hline
v1.0.0(14) & $2023-12-06$ & \quantities{Bresolin G.} & Orlandi G. & Aggiornamento UC e requisiti presenti. Introduzione diagrammi dei casi d'uso. Aggiunta di nuovi UC e nuovi requisiti.\\
\hline
v1.0.0(5) & $2023-11-23$ & \quantities{Ciriolo I.} & Campese M. & Introduzione, Descrizione, Studio dei primi Casi d'uso, Studio dei primi requisiti.\\

    
\end{xltabular}}
\newpage

% INDICE
\tableofcontents
\newpage

% INTRODUZIONE
\section{Introduzione}
\subsection{Scopo del documento}
Il documento ha lo scopo di definire i casi d'uso della \textit{web application} che verrà sviluppata nel corso del progetto. Vengono inoltre presentati i requisiti funzionali, di qualità e di vincolo. %e le funzionalità aggiuntive che saranno implementate non essendo ritenute essenziali/necessarie.
\\
%\subsection{Scopo del prodotto}
\subsection{Glossario}
%\paragraph{}Per evitare ambiguità e incomprensioni relative al linguaggio e ai termini utilizzati nella documentazione relativa al progetto viene presentato un Glossario. I termini ambigui o specifici presenti nello stesso, verranno identificati con un pedice |g|. (DA VALUTARE COME RICONOSCERE IL TERMINE )

\subsection{Riferimenti}
   \subsubsection{Riferimenti normativi}
   \begin{itemize}
    \item \textit{Regolamento del progetto didattico}: \\
    \href{https://www.math.unipd.it/~tullio/IS-1/2023/Dispense/PD2.pdf}{https://www.math.unipd.it/\textasciitilde tullio/IS-1/2023/Dispense/PD2.pdf};
    \item \textit{\S Norme di progetto}.
    \end{itemize}
    
    \subsubsection{Riferimenti informativi}
    
    \begin{itemize}
   % \item \textit{Glossario v0.0.1} (da creare parallelamente); 
    \item \textit{Presentazione capitolato C1}:\\
    \href{https://www.math.unipd.it/~tullio/IS-1/2023/Progetto/C1.pdf}{https://www.math.unipd.it/\textasciitilde tullio/IS-1/2023/Progetto/C1.pdf};
    \item \textit{Verbali esterni ed interni};
    \item \textit{Analisi e descrizione delle funzionalità: Use Case e relativi diagrammi (UML):}\\
    \href{https://www.math.unipd.it/~rcardin/swea/2022/Diagrammi%20Use%20Case.pdf}{https://www.math.unipd.it/\textasciitilde rcardin/swea/2022/Diagrammi\%20Use\%20Case.pdf ;}
    \item \textit{Analisi e descrizione delle funzionalità: Diagrammi delle Attività (UML):}\\
    \href{https://www.math.unipd.it/~rcardin/swea/2022/Diagrammi%20di%20Attivit%C3%A0.pdf}{https://www.math.unipd.it/\textasciitilde rcardin/swea/2022/Diagrammi\%20di\%20Attivit\%C3\%A0.pdf;}
    

    \item \textit{Progettazione e Programmazione: Diagrammi delle Classi:}\\
    \href{https://www.math.unipd.it/~rcardin/swea/2023/Diagrammi%20delle%20Classi.pdf}{https://www.math.unipd.it/\textasciitilde rcardin/swea/2023/Diagrammi\%20delle\%20Classi.pdf.}
    \end{itemize}

% DESCRIZIONE
\newpage
\section{Descrizione}
\subsection{Obiettivi del prodotto}
Il proponente richiede la creazione di un sito web dotato di un chatbot realizzato tramite interrogazioni ad un modello di intelligenza artificiale, che sia d'aiuto a tutto il personale dell'azienda.
In particolare il sistema deve facilitare l'accesso dei dipendenti alle informazioni, regolamentazioni e alle indicazioni aziendali.\\
Il prodotto deve fornire nuove modalità di formazione e assistenza sul lavoro, superando i limiti di un classico \textit{repository} documentale, in maniera tale da rendere le informazioni più fruibili, concentrandosi più sui contenuti; il sistema deve riuscire a comunicare con i lavoratori in modo semplice, in maniera tale che qualsiasi persona, anche la meno esperta, riesca a ricavare le informazioni che cerca.
%Infine, deve riuscire a migliorare il comportamento dei lavoratori, facendo in modo che rispettino le regole.
\subsection{Funzioni del prodotto}
Il software deve essere in grado di ricevere in input documenti, in particolare \textit{PDF}, oppure documenti Word, con i quali possa rispondere a qualsiasi domanda sugli stessi sottoposta dall'utente. 
\subsection{Caratteristiche del prodotto}
Il sistema deve poter aver accesso a tutti i file caricati, non può rispondere simultaneamente a più utenti ma ne può gestirne uno alla volta e deve poter riconoscere la lingua italiana e quella inglese. I documenti caricati devono poter rimanere accessibili così come lo storico delle richieste effettuate dall'utente.
\subsection{Utenti}
Gli utenti possono interagire con il sistema attraverso un'interfaccia dove è presente un \textit{chatbot} ma possono anche visualizzare i file che il sistema gestisce. L'utente può chiedere informazioni riguardo ai documenti specifici attraverso la risposta data dal \textit{chatbot}. Se l'utente chiede al sistema informazioni non pertinenti, il sistema deve notificarlo attraverso una risposta cortese.
% CASI D'USO
\newpage
\section{Casi d'uso}

\subsection{Obiettivi}
La seguente sezione ha l'obiettivo di presentare i diversi Casi d'Uso relativi all'applicazione \textit{Knowledge Management AI}. 
\subsection{Notazione} 
I Casi d'Uso sono indicati con la notazione seguente: \\ UC[X].[Y] \ \  in cui:
\begin{itemize}
\item UC sta per \textit{Use Case};
\item .[X] indica il numero del Caso d'Uso, presentati secondo successione gerarchica;
\item .[Y] indica il possibile numero del sotto-Caso d'uso, che aggiunge informazioni al Caso d'Uso [X]. Anche i sotto-Casi d'Uso sono presentati secondo successione gerarchica.
\end{itemize}

\subsection{Attori}
L'applicazione \textit{Knowledge Management AI} richiede un singolo attore che rappresenta l'utente (ovvero il dipendente aziendale) che interagisce con il ChatBot. \\ Si specifica che l'azienda proponente non richiede l'autenticazione degli utenti nell'applicazione.%, poiché rappresenterebbe un rallentamento alla divulgazione rapida e semplice di documenti all'interno della rete aziendale; Per questo motivo l'autenticazione non viene presentata tra i requisiti.
\subsection{Didascalie}
Per ogni Caso d'Uso è presentata una breve didascalia che ne indica:
\begin{itemize}
    \item Descrizione;
    \item Scenario (ovvero il modo in cui l'attore giunge al Caso d'Uso tramite passi indicati da elenco numerato);
    \item Estensioni;
    \item Precondizioni (ovvero ciò che accade prima che il Caso d'Uso si presenti);
    \item Postcondizioni (ovvero ciò che accade dopo che il Caso d'Uso si è presentato).
\end{itemize}
N.B: non viene specificato l'attore principale di ogni Caso d'Uso poiché è stato precedentemente specificato che si tratta sempre di un singolo utente che si interfaccia con l'applicazione.
\subsection{Lista dei casi d'uso}
\begin{itemize}

    \item UC1 - Visualizzazione interfaccia Chatbot;
    \begin{itemize}
        \item \textbf{Descrizione:} l'utente desidera interagire con il chatbot;
        \item \textbf{Scenario: } 
        \begin{enumerate}[label={\arabic*.}, align=left]
        \item l'utente apre l'applicazione;
        \end{enumerate}
        \item \textbf{Estensioni: };
        \item \textbf{Precondizioni:} l'utente apre l'applicazione;
        \item \textbf{Postcondizioni:} l'utente interagisce con l'interfaccia chatbot.\\
    \end{itemize}
    
    \item UC2 - Visualizzazione documenti;
    \begin{itemize}
        \item \textbf{Descrizione:} l'utente visualizza la lista di documenti presenti nell'applicazione;
        \item \textbf{Scenario: }
        \begin{enumerate}[label={\arabic*.}, align=left]
        \item l'utente apre l'applicazione;
        \item l'utente visualizza una lista di documenti.
        \end{enumerate}
        \item \textbf{Estensioni: };
        \item \textbf{Precondizioni:} l'utente apre l'applicazione;
        \item \textbf{Postcondizioni:} l'utente visualizza i documenti che si trovano nell'applicazione.\\
    \end{itemize}
    
    \item UC3 - Selezione documento singolo;
    \begin{itemize}
        \item \textbf{Descrizione: }l'utente seleziona un documento presente nell'applicazione;
        \item \textbf{Scenario: }
        \begin{enumerate}[label={\arabic*.}, align=left]
        \item l'utente apre l'applicazione;
        \item l'utente sfoglia la lista di documenti;
        \item l'utente seleziona uno dei documenti.
        \end{enumerate}
        \item \textbf{Estensioni: };
        \item \textbf{Precondizioni:} l'utente apre l'applicazione e sceglie un documento;
        \item \textbf{Postcondizioni:} l'utente ha selezionato il documento desiderato dalla lista.\\
    \end{itemize}
    
    \item UC4 - Selezione insieme di documenti;
    \begin{itemize}
        \item \textbf{Descrizione: }l'utente seleziona più di un documento;
        \item \textbf{Scenario: }
        \begin{enumerate}[label={\arabic*.}, align=left]
        \item l'utente entra nell'applicazione;
        \item l'utente sfoglia la lista dei documenti;
        \item l'utente sceglie una serie di documenti da selezionare.
        \end{enumerate}
        \item \textbf{Estensioni: };
        \item \textbf{Precondizioni:} l'utente apre l'applicazione e sceglie una serie di documenti dalla lista;
        \item \textbf{Postcondizioni:} l'utente ha selezionato i documenti desiderati dalla lista.\\
    \end{itemize}
    
    \item UC5 - Caricamento singolo documento;
    \begin{itemize}
        \item \textbf{Descrizione: }l'utente carica un documento nell'applicazione;
        \item \textbf{Scenario: }
        \begin{enumerate}[label={\arabic*.}, align=left]
        \item l'utente apre l'applicazione;
        \item l'utente seleziona un file da caricare;
        \item l'utente carica un documento.
        \end{enumerate}
        \item \textbf{Estensioni: }UC5.1;
        \item \textbf{Precondizioni:} l'utente apre l'applicazione;
        \item \textbf{Postcondizioni:} l'utente vede il completamento del caricamento del singolo file.
    \end{itemize}
    \item UC5.1 - Caricamento di più documenti;
    \begin{itemize}
        \item \textbf{Descrizione: }l'utente carica più documenti nell'applicazione;
        \item \textbf{Scenario: }
        \begin{enumerate}[label={\arabic*.}, align=left]
        \item l'utente apre l'applicazione;
        \item l'utente seleziona più file da caricare;
        \item l'utente carica più documenti.
        \end{enumerate}
        \item \textbf{Estensioni: };
        \item \textbf{Precondizioni:} l'utente apre l'applicazione;
        \item \textbf{Postcondizioni:} l'utente vede il completamento del caricamento massivo.\\
    \end{itemize}
    
    \item UC6 - Visualizzazione completamento caricamento;
    \begin{itemize}
        \item \textbf{Descrizione: }l'utente vede che il caricamento del suo file (o dei suoi file) è andato a buon fine;
        \item \textbf{Scenario: }
        \begin{enumerate}[label={\arabic*.}, align=left]
        \item l'utente apre l'applicazione;
        \item l'utente seleziona uno o più file da caricare;
        \item l'utente carica uno o più documenti;
        \item il caricamento va a buon fine.
        \end{enumerate}
        \item \textbf{Estensioni: };
        \item \textbf{Precondizioni:} l'utente carica uno o più file (UC5, UC5.1);
        \item \textbf{Postcondizioni:} l'utente verifica che il caricamento sia andato a buon fine.\\
    \end{itemize}
    
    \item UC7 - Visualizzazione di errore nel caricamento;
    \begin{itemize}
        \item \textbf{Descrizione: }l'utente vede che è avvenuto un errore nel caricamento del file (o dei file);
        \item \textbf{Scenario: }
        \begin{enumerate}[label={\arabic*.}, align=left]
        \item l'utente apre l'applicazione;
        \item l'utente seleziona uno o più file da caricare;
        \item l'utente carica uno o più documenti;
        \item il caricamento non va a buon fine.
        \end{enumerate}
        \item \textbf{Estensioni: }UC7.1, \ UC7.2, \ UC7.3;
        \item \textbf{Precondizioni:} caricamento di uno o più documenti (UC5, UC5.1);
        \item \textbf{Postcondizioni:} l'utente vede l'avviso di errore nel caricamento.
    \end{itemize}
    
    \item UC7.1 - Visualizzazione di errore: documento già presente;   
    \begin{itemize}
    \item \textbf{Descrizione: }l'utente vede che è avvenuto un errore nel caricamento del file (o dei file) poiché il documento caricato è già stato inserito nell'applicazione;
        \item \textbf{Scenario: }
        \begin{enumerate}[label={\arabic*.}, align=left]
        \item scenario di caricamento documento/i UC5 o UC5.1
        \item l'utente carica nuovamente un documento caricato in precedenza;
        \item il caricamento non va a buon fine.
        \end{enumerate}
        \item \textbf{Estensioni: };
        \item \textbf{Precondizioni:} caricamento di uno o più documenti (UC5, UC5.1) e caricamento di un documento già caricato;
        \item \textbf{Postcondizioni:} l'utente vede l'avviso di errore nel caricamento in cui viene scritto che il documento è già presente nell'applicazione.
    \end{itemize}

    \item UC7.2 - Visualizzazione di errore: documento vuoto;
    \begin{itemize}
    \item \textbf{Descrizione: }l'utente vede che è avvenuto un errore nel caricamento del file (o dei file) poichè il documento caricato è privo di contenuto;
        \item \textbf{Scenario: }
        \begin{enumerate}[label={\arabic*.}, align=left]
        \item l'utente apre l'applicazione;
        \item l'utente seleziona uno o più file da caricare;
        \item l'utente carica uno o più documenti vuoti;
        \item il caricamento non va a buon fine.
        \end{enumerate}
        \item \textbf{Estensioni: };
        \item \textbf{Precondizioni:} caricamento di uno o più documenti (UC5, UC5.1) e caricamento di almeno un documento vuoto;
        \item \textbf{Postcondizioni:} l'utente vede l'avviso di errore nel caricamento in cui viene scritto che il documento inserito è vuoto.
    \end{itemize}

    \item UC7.3 - Visualizzazione di errore: formato del documento non consentito;
    \begin{itemize}
    \item \textbf{Descrizione: }l'utente vede che è avvenuto un errore nel caricamento del file (o dei file) poichè il formato del documento caricato non è riconoscibile dall'applicazione;
        \item \textbf{Scenario: }
        \begin{enumerate}[label={\arabic*.}, align=left]
        \item scenario di caricamento documento/i UC5 o UC5.1;
        \item il caricamento non va a buon fine.
        \end{enumerate}
        \item \textbf{Estensioni: };
        \item \textbf{Precondizioni :}caricamento di uno o più documenti (UC5, UC5.1);
        \item \textbf{Postcondizioni:} l'utente vede l'avviso di errore nel caricamento in cui viene scritto che il formato del documento non è compatibile con l'applicazione.\\
    \end{itemize}
       
    \item UC8 - Cancellazione singolo documento;
    \begin{itemize}
        \item \textbf{Descrizione: }l'utente elimina un documento presente all'interno dell'applicazione;
        \item \textbf{Scenario: }
        \begin{enumerate}[label={\arabic*.}, align=left]
        \item l'utente seleziona un documento (UC3);
        \item l'utente elimina un documento.
        \end{enumerate}
        \item \textbf{Estensioni: };
        \item \textbf{Precondizioni :}nell'applicazione è presente almeno un documento;
        \item \textbf{Postcondizioni: }il modello di IA viene addestrato nuovamente senza il documento cancellato.\\
    \end{itemize}

    \item UC9 - Cancellazione insieme di documenti;
    \begin{itemize}
        \item \textbf{Descrizione: }l'utente elimina più di un documento presente all'interno dell'applicazione;
        \item \textbf{Scenario: }
        \begin{enumerate}[label={\arabic*.}, align=left]
        \item l'utente seleziona più documenti (UC4);
        \item l'utente elimina i documenti.
        \end{enumerate}
        \item \textbf{Estensioni: };
        \item \textbf{Precondizioni :}nell'applicazione è presente almeno un documento;
        \item \textbf{Postcondizioni: }il modello di IA viene addestrato nuovamente senza i documenti cancellati.\\ 
    \end{itemize}
  
    \item UC10 - Visualizzazione di domande pre-impostate;
    \begin{itemize}
        \item \textbf{Descrizione: }l'utente visualizza delle domande predefinite con le quali può interagire con l'interfaccia chatbot;
        \item \textbf{Scenario: }
        \begin{enumerate}[label={\arabic*.}, align=left]
            
        \item l'utente apre l'applicazione;
        \item l'utente visualizza le domande predefinite;
        \end{enumerate}
        \item \textbf{Estensioni: };
        \item \textbf{Precondizioni:} l'utente apre l'applicazione;
        \item \textbf{Postcondizioni:} l'utente seleziona una domanda.\\
    \end{itemize}

    \item UC11 - Selezione domanda pre-impostata;
    \begin{itemize}
        \item \textbf{Descrizione: }l'utente seleziona una domanda predefinita con la quale può interagire con l'interfaccia chatbot;
        \item \textbf{Scenario: }
        \begin{enumerate}[label={\arabic*.}, align=left]
        \item l'utente apre l'applicazione;
        \item l'utente visualizza le domande predefinite;
        \item l'utente seleziona una delle domande proposte.
        \end{enumerate}
        \item \textbf{Estensioni: };
        \item \textbf{Precondizioni:} l'utente apre l'applicazione;
        \item \textbf{Postcondizioni:} l'utente visualizza la risposta alla domanda selezionata.\\
    \end{itemize}
    
    \item UC12 - Inserimento domanda o testo;
    \begin{itemize}
        \item \textbf{Descrizione: }l'utente pone un quesito al chatbot;
        \item \textbf{Scenario: }
        \begin{enumerate}[label={\arabic*.}, align=left]
        \item l'utente apre l'applicazione;
        \item l'utente inserisce una domanda per il chatbot;
        \item l'utente invia la domanda.
        \end{enumerate}
        \item \textbf{Estensioni: }UC13, UC14;
        \item \textbf{Precondizioni:} l'utente apre l'applicazione;
        \item \textbf{Postcondizioni: }l'utente invia il quesito e visualizza la risposta alla domanda posta (o un messaggio alternativo).\\
    \end{itemize}

    \item UC13 - Digitazione domanda o testo;
    \begin{itemize}
        \item \textbf{Descrizione: }l'utente pone un quesito al chatbot digitando la domanda sulla tastiera;
        \item \textbf{Scenario: }
        \begin{enumerate}[label={\arabic*.}, align=left]
        \item l'utente apre l'applicazione;
        \item l'utente digita una domanda per il chatbot;
        \item l'utente invia la domanda.
        \end{enumerate}
        \item \textbf{Estensioni: };
        \item \textbf{Precondizioni:} l'utente apre l'applicazione;
        \item \textbf{Postcondizioni: }l'utente visualizza la risposta alla domanda posta.\\
    \end{itemize}
        
        \item UC14 - Inserimento vocale domanda o testo;
        \begin{itemize}
        \item \textbf{Descrizione: }l'utente pone un quesito al chatbot tramite microfono;
        \item \textbf{Scenario: }
        \begin{enumerate}[label={\arabic*.}, align=left]
        \item l'utente apre l'applicazione;
        \item l'utente seleziona l'opzione di inserimento vocale (UC21);
        \item l'utente pone vocalmente la domanda;
        \item l'utente invia la domanda.
        \end{enumerate}
        \item \textbf{Estensioni: };
        \item \textbf{Precondizioni:} l'utente seleziona l'opzione di inserimento vocale (UC21);
        \item \textbf{Postcondizioni: }l'utente visualizza la risposta alla domanda posta.\\
    \end{itemize}
    
    \item UC15 - Invio domanda o testo;
    \begin{itemize}
        \item \textbf{Descrizione: }l'utente dopo aver inserito una domanda o un testo, lo invia al chatbot;
        \item \textbf{Scenario: }
        \begin{enumerate}[label={\arabic*.}, align=left]
        \item l'utente apre l'applicazione;
        \item l'utente pone un quesito al chatbot (UC12);
        \item l'utente invia la domanda.
        \end{enumerate}
        \item \textbf{Estensioni: }UC16;
        \item \textbf{Precondizioni:} l'utente inserisce un quesito per il chatbot (UC12);
        \item \textbf{Postcondizioni: }l'utente visualizza la risposta alla domanda posta.\\
    \end{itemize}
    
    \item UC16 - Visualizzazione risposta cortese;
    \begin{itemize}
        \item \textbf{Descrizione: }l'utente visualizza una risposta cortese dopo aver posto un quesito ambiguo all'applicazione;
        \item \textbf{Scenario: }
        \begin{enumerate}[label={\arabic*.}, align=left]
        \item l'utente apre l'applicazione;
        \item l'utente pone un quesito ambiguo al chatbot;
        \item l'utente invia la domanda.
        \end{enumerate}
        \item \textbf{Estensioni: };
        \item \textbf{Precondizioni:} l'utente invia un quesito (UC13);
        \item \textbf{Postcondizioni: }l'utente pone un ulteriore quesito al chatbot.\\
    \end{itemize}
    
    \item UC17 - Visualizzazione risposta;
    \begin{itemize}
        \item \textbf{Descrizione: }l'utente visualizza una risposta alla domanda che ha posto (o che ha selezionato);
        \item \textbf{Scenario: }
        \begin{enumerate}[label={\arabic*.}, align=left]
        \item l'utente apre l'applicazione;
        \item l'utente inserisce un quesito;
        \item l'utente invia il quesito;
        \item l'utente riceve la risposta a tale quesito.
        \end{enumerate}
        \item \textbf{Estensioni: };
        \item \textbf{Precondizioni:} l'utente invia un quesito (UC13);
        \item \textbf{Postcondizioni: }l'utente prosegue l'interazione con il chatbot.\\
    \end{itemize}
    
    \item UC18 - Visualizzazione storico chat;
    \begin{itemize}
        \item \textbf{Descrizione: }l'utente visualizza lo storico delle chat nell'applicazione;
        \item \textbf{Scenario: }
        \begin{enumerate}[label={\arabic*.}, align=left]        
        \item l'utente apre l'applicazione;
        \item l'utente visualizza lo storico delle chat;
        \end{enumerate}
        \item \textbf{Estensioni: } UC18.1, UC18.1.1, UC18.1.2;
        \item \textbf{Precondizioni:} l'utente apre l'applicazione;
        \item \textbf{Postcondizioni: }l'utente può leggere chat salvate in precedenza.
    \end{itemize}
    
    \item UC18.1 - Visualizzazione singola chat;
    \begin{itemize}
        \item \textbf{Descrizione: }l'utente visualizza una chat dallo storico delle chat nell'applicazione;
        \item \textbf{Scenario: }
        \begin{enumerate}[label={\arabic*.}, align=left]
        \item l'utente apre l'applicazione;
        \item l'utente visualizza lo storico delle chat;
        \item l'utente seleziona una chat.
        \end{enumerate}
        \item \textbf{Estensioni: }UC18.1.1, UC18.1.2;
        \item \textbf{Precondizioni:} l'utente visualizza lo storico delle conversazioni;
        \item \textbf{Postcondizioni: }l'utente legge una conversazione salvata nello storico.
    \end{itemize}
    
    \item UC18.1.1 - Visualizzazione titolo chat;
    \begin{itemize}
         \item \textbf{Descrizione: }l'utente visualizza il titolo di una conversazione dallo storico delle chat nell'applicazione;
        \item \textbf{Scenario: }
        \begin{enumerate}[label={\arabic*.}, align=left]
        \item l'utente apre l'applicazione;
        \item l'utente visualizza lo storico delle chat;
        \item l'utente seleziona una chat;
        \item l'utente ne visualizza il titolo.
        \end{enumerate}
        \item \textbf{Estensioni: };
        \item \textbf{Precondizioni:} l'utente visualizza una chat dallo storico delle conversazioni;
        \item \textbf{Postcondizioni: }l'utente legge il titolo della conversazione salvata nello storico.
    \end{itemize}
    
    \item UC18.1.2 - Visualizzazione orario chat;
    \begin{itemize}
        \item \textbf{Descrizione: }l'utente visualizza l'orario nel quale una conversazione dallo storico delle chat è avvenuta;
        \item \textbf{Scenario: }
        \begin{enumerate}[label={\arabic*.}, align=left]
        \item l'utente apre l'applicazione;
        \item l'utente visualizza lo storico delle chat;
        \item l'utente seleziona una chat;
        \item l'utente ne visualizza l'orario.
        \end{enumerate}
        \item \textbf{Estensioni: };
        \item \textbf{Precondizioni:} l'utente visualizza una chat dallo storico delle conversazioni;
        \item \textbf{Postcondizioni: }l'utente può leggere l'orario della conversazione salvata nello storico.\\
    \end{itemize}

    \item UC19 - Salvataggio conversazione;
    \begin{itemize}
        \item \textbf{Descrizione: }l'utente salva una serie di scambi di domande e risposte avvenute con il chatbot;
        \item \textbf{Scenario: }
        \begin{enumerate}[label={\arabic*.}, align=left]
        \item l'utente apre l'applicazione;
        \item l'utente inserisce ed invia una serie di quesiti al chatbot (UC13), ricevendo una serie di risposte (UC15);
        \item l'utente salva lo scambio di domande e riposte.
        \end{enumerate}
        \item \textbf{Estensioni: };
        \item \textbf{Precondizioni:} l'utente invia dei quesiti (UC13);
        \item \textbf{Postcondizioni: }l'utente potrà accedere alla conversazione in qualsiasi momento.\\
    \end{itemize}

     \item UC20 - Eliminazione conversazione;
    \begin{itemize}
        \item \textbf{Descrizione: }l'utente elimina una serie di scambi di domande e risposte avvenute con il chatbot;
        \item \textbf{Scenario: }
        \begin{enumerate}[label={\arabic*.}, align=left]
        \item l'utente apre l'applicazione;
        \item l'utente inserisce ed invia una serie di quesiti al chatbot (UC13), ricevendo una serie di risposte (UC15);
        \item l'utente elimina lo scambio di domande e riposte.
        \end{enumerate}
        \item \textbf{Estensioni: };
        \item \textbf{Precondizioni:} l'utente invia dei quesiti (UC 13);
        \item \textbf{Postcondizioni: }il modello di IA viene addestrato nuovamente.\\
    \end{itemize}

    \item UC21 - Selezione inserimento vocale;
    \begin{itemize}
        \item \textbf{Descrizione: }l'utente seleziona l'opzione di inserimento vocale di quesiti;
        \item \textbf{Scenario: }
        \begin{enumerate}[label={\arabic*.}, align=left]
        \item l'utente apre l'applicazione;
        \item l'utente seleziona l'opzione di inserimento vocale.
        \end{enumerate}
        \item \textbf{Estensioni: };
        \item \textbf{Precondizioni:} l'utente apre l'applicazione;
        \item \textbf{Postcondizioni: }l'utente inserisce vocalmente un quesito.\\
    \end{itemize}

    
\end{itemize}

%\renewcommand{\arraystretch}{1.5}
%\begin{tabularx}{\textwidth}{|c|X|c|c|c|}\hline
%\textbf{Caso d'uso} & \textbf{Attori} & \textbf{Precondizioni} & %\textbf{Postcondizioni} & \textbf{Scenario principale} \\
%\end{tabularx}


% REQUISITI
\newpage
\section{Requisiti}
Ogni requisito preso in analisi sarà trattato con la sigla:\\
R[Tipo].[Importanza].[Ordine];\\
nella quale:\\
\begin{itemize}
\item R indica la parola "requisito";
\item Tipo può essere:
    \begin{itemize}
	\item F (Funzionale);
	\item Q (Qualità);
	\item V (Vincolo).
    \end{itemize}
\item Importanza può essere:
    \begin{itemize}
        \item O (Obbligatorio);
	\item D (Desiderabile);
	\item OP (Opzionale).
    \end{itemize}
\item Ordine: indica l’ordine in cui il requisito è stato preso in esame. Per i requisiti con stessa importanza e tipologia indica l’ordine cronologico di precedenza che esso ha avuto rispetto agli altri, ovvero il suo ordine inserimento nell'Analisi dei requisiti.
\end{itemize}
N.B: per requisiti che dovessero generare requisiti figli, si andrà ad aggiungere .[N.figlio] alla notazione successivamente all’ordine, di modo da avere i figli ordinati per importanza.

\subsection{Requisiti funzionali}
Di seguito la specifica per i requisiti funzionali, i quali descrivono le funzionalità del sistema, le azioni
che il sistema può compiere e le informazioni che il sistema può fornire.
Si intende che la numerologia di ciascuno rispecchi ove previsto gli UC presenti, mentre le sigle sotto
riportate indicano:
\begin{itemize}
    \item ROF - Requisito Obbligatorio Funzionale;
    \item RDF - Requisito Desiderabile Funzionale;
    \item ROPF - Requisito Opzionale Funzionale.
\end{itemize}

\renewcommand{\arraystretch}{1.5}
\begin{tabularx}{\textwidth}{|c|X|c|c}\hline
\textbf{Codice} & \textbf{Descrizione}  & \textbf{Fonti} \\
\hline ROF1\ & Il lavoratore deve poter interagire con il chatbot tramite un'interfaccia semplice. &  UC1 \\
\hline ROF2\ & L'utente deve poter caricare uno o più documenti & UC5, UC5.1, UC6 \\
\hline ROF3\ & L'utente deve poter eliminare uno o più documenti & UC8, UC9 \\
\hline ROF4\ & Il sistema deve visualizzare un messaggio di errore esplicativo &  UC7, UC7.1, UC7.2, UC7.3 \\
\hline ROF5\ & Il sistema deve poter memorizzare lo storico delle chat & UC18, UC19 \\
\hline ROPF6\ & L’utente deve poter visualizzare le lingue secondarie presenti e i dettagli della singola lingua &  UC14\\
\hline ROF7\ & Il sistema può interagire con l'utente tramite comando vocale & UC14, UC21 \\
\hline ROF8 \ & L'utente deve poter visualizzare domande preimpostate & UC11  \\
\hline ROF9 \ & L'utente deve poter scrivere domande & UC12, UC13  \\
%\hline RDF9 \ & Il sistema deve poter memorizzare la conversazione avuta con il lavoratore & UC16 \\
\hline
\end{tabularx}

\subsection{Requisiti di qualità}
Di seguito la specifica per i requisiti di qualità, i quali descrivono come un sistema deve essere, o
come il sistema deve esibirsi, per soddisfare le esigenze dell'utente.
Le sigle sotto riportate possono essere così classificate:
\begin{itemize}
    \item ROQ - Requisito Obbligatorio di Qualità;
    \item RDQ - Requisito Desiderabile di Qualità;
    \item ROPQ - Requisito Opzionale di Qualità.
\end{itemize}
\renewcommand{\arraystretch}{1.5}
\begin{tabularx}{\textwidth}{|c|X|c|c}\hline
\textbf{Codice} & \textbf{Descrizione}  & \textbf{Fonti} \\
\hline
ROQ1 & Il sistema deve garantire il servizio per almeno un utente & - \\
\hline
% \hline ROQ 1\ & Manuale utente disponibile all’utilizzo &  Norme di Progetto \\
% \hline ROPQ 2\ & Documentazione e manuale sviluppatore disponibile alla consultazione & Norme di Progetto \\
% \hline ROQ 3\ & Il codice sorgente deve essere presente su GitHub  &  Capitolato \\
% \hline ROPQ 4\ & Documentazione dettagliata delle tecnologie usate & Capitolato \\
% \hline
\end{tabularx}

\subsection{Requisiti di vincolo}
Di seguito la specifica per i requisiti di vincolo, i quali descrivono i limiti e le restrizioni normative/legislative che un sistema
deve rispettare per soddisfare le esigenze dell'utente.
Le sigle sotto riportate possono essere così classificate:
\begin{itemize}
    \item ROV - Requisito Obbligatorio di Vincolo;
    \item RDV - Requisito Desiderabile di Vincolo;
    \item ROPV - Requisito Opzionale di Vincolo. 
\end{itemize}
\renewcommand{\arraystretch}{1.5}
\begin{tabularx}{\textwidth}{|c|X|c|c}\hline
\textbf{Codice} & \textbf{Descrizione} & \textbf{Fonti} \\
\hline
ROV1 & Il sistema deve garantire la privacy delle informazioni sensibili condivise dall'utente per l'addestramento del modello. & - \\
% \hline ROV 1 \ & Utilizzo di tecnologie con intelligenza artificiale & Capitolato  \\
% \hline ROV 2 \ & Algoritmi avanzati di ricerca e comprensione del linguaggio naturale &  Capitolato \\
% \hline ROV 3 \ & Utilizzo di componenti come \textit{OpenAI} , come motore per le funzionalità di \textit{NPL}(comprensione del testo) & Capitolato \\
% \hline ROPV 4 \ & Utilizzo di \textit{LangChain} (progetto opensource che permette di integrare modelli di AI) & Capitolato \\

\hline

\end{tabularx}




\end{document}
